\documentclass{article}
\begin{document}
\section*{All-length ordered coefficient fold}
Fix a list of occurrence addresses. The coefficient objects and arrows
vary pointwise in the category of functors from the discrete address
type. Distinct occurrences may vary independently when labelled by
distinct addresses, even if their carrier values coincide. Repeating a
label in the input list repeats the same coefficient, rather than
inventing a different address. An arbitrary finite ordered family is
covered by choosing its positions as addresses.

The functor takes the right-associated, unit-terminated tensor fold.
Its cotensorator is the existing occurrence shuffle, not a new
coherence witness of unspecified construction. Its counit collapses the
fold of units, recursively through actual unitors, with identity at
length zero. The two iterated shuffles factor through independent
recursive triple folds; associativity is derived by induction from the
local middle-four associator and the nullary unit law. The two unit
diagrams are likewise derived, and all three laws instantiate with
the actual braided tensor and unitors. Invertible counit and
cotensorator give a strong monoidal functor in Mathlib.

This establishes naturality for pointwise coefficient morphisms and
the ordinary monoidal associativity and left/right unit equations for
one ordered presentation. Separately, in a symmetric monoidal
structure a constructive deletion normal form proves equality of
all typed adjacent-swap circuits with the same source and target words
of distinct tags. Coefficient labels may repeat. The local deletion
square uses Yang--Baxter under braiding and requires symmetry for
involutivity. This additive circuit result does not itself establish
monoidality of the reindexing maps for multiplicative ordered folds.

The standard symmetric additive whole-fibre direct image has explicit
object, arrow, singleton, and split constructions; all naturality,
units, and the native split pentagon are proved. A permutation of
outer row tags lifts to a typed circuit on flattened descendants.
A padded-tail proof identifies the lift with the actual block
braiding. A row can have no descendants without dropping its unit
in the row chart. Both native pentagon paths are independently
identified with one final ordered iterate of sums following tagged
circuits, which are equal by the deletion normal form. This yields a
native \texttt{Push}, pinned to the addressed finite sums; singleton
and split maps are isomorphisms. As a small discriminator, a
noninjective map from two positions to two positions sends two equal
coefficient objects to an ordered two-term sum and has an actual
empty second fibre.

This supplies the additive finite direct-image input for native
families and augmented-family morphisms. It does not identify the
pinned raw-chart comparisons with all original dependent and
independent morphisms, prove their naturality in carrier-changing
maps at arbitrary arity, or prove all-arity dependent augmented
Fubini. The binary native augmentation theorem remains the current
Fubini boundary. Neither initiality of zero nor a choice of
elementwise restrictions of whole-fibre arrows enters this result.
\end{document}
